\chapter{30 天学习与长期维护路线}
\label{chap:roadmap}

\begin{takeaway}[学习目标]
把全书转成每日可验收任务，建立“读原始资料、写最小实现、跑评测、复盘失败”的学习循环，并为长期跟踪高速变化的 Agent 生态建立维护机制。
\end{takeaway}

\section{每天必须留下一个产物}

学习 Agent 最容易变成收藏链接。每天至少留下代码、测试、评测样本、架构图、trace、技术笔记或演示中的一种。输入资料的时间不应超过总学习时间的一半。

\section{第 1 周：建立最小闭环}

\begin{longtable}{>{\bfseries}p{0.12\textwidth} p{0.38\textwidth} p{0.4\textwidth}}
\toprule
天 & 学习与实践 & 当日产物 \\
\midrule
1 & 区分调用、工作流与 Agent & 一个场景的复杂度评分和任务合同 \\
2 & 学习六层结构 & Agent 数据流与三个故障归层 \\
3 & 阅读有效 Agent 模式 & 同一任务的三种控制流草图 \\
4 & 运行最小 Agent & 可执行循环与结构化轨迹 \\
5 & 接入一个真实模型 & 模型适配层和 10 条协议测试 \\
6 & 加入停止与预算 & 五种停止原因及测试 \\
7 & 周复盘 & Demo、失败清单、下周基线 \\
\bottomrule
\end{longtable}

\section{第 2 周：工具、上下文与知识}

\begin{longtable}{>{\bfseries}p{0.12\textwidth} p{0.38\textwidth} p{0.4\textwidth}}
\toprule
天 & 学习与实践 & 当日产物 \\
\midrule
8 & 写三个工具合同 & schema、错误码、副作用分级 \\
9 & 实现权限与 dry-run & proposal/commit 演示 \\
10 & 实现幂等与重试 & 重复请求不会重复写入的测试 \\
11 & 设计上下文预算 & 组装器与压缩策略 \\
12 & 建立小型文档索引 & 20--50 篇文档及元数据 \\
13 & 做检索与引用评测 & 30 题黄金集和 Recall@5 \\
14 & 周复盘 & 比较三种检索方案并写结论 \\
\bottomrule
\end{longtable}

\section{第 3 周：编排、评测与安全}

\begin{longtable}{>{\bfseries}p{0.12\textwidth} p{0.38\textwidth} p{0.4\textwidth}}
\toprule
天 & 学习与实践 & 当日产物 \\
\midrule
15 & 建立结构化计划 & 含依赖、预算和验收的计划 \\
16 & 实现 worker 委托 & 委托合同与上下文裁剪 \\
17 & 并行与合并 & 冲突识别和去重测试 \\
18 & 建立 50 条评测集 & 分层标签和确定性断言 \\
19 & 加 trace 与指标 & 成功和失败时间线 \\
20 & 画威胁模型 & 资产、主体、边界和缓解措施 \\
21 & 红队与复盘 & 20 条攻击样本及修复回归 \\
\bottomrule
\end{longtable}

\section{第 4 周：做成可交付产品}

\begin{longtable}{>{\bfseries}p{0.12\textwidth} p{0.38\textwidth} p{0.4\textwidth}}
\toprule
天 & 学习与实践 & 当日产物 \\
\midrule
22 & 确定毕业项目范围 & 一页立项书与非目标 \\
23 & 完成端到端 happy path & 一条真实任务跑通 \\
24 & 加检查点与恢复 & 故障注入演示 \\
25 & 设计确认和失败界面 & 可点击原型或 CLI 交互 \\
26 & 跑完整评测 & 基线报告和所有回归 \\
27 & 用户测试 & 五位用户的观察与问题列表 \\
28 & 性能和成本优化 & 前后对比表 \\
29 & 整理开源仓库 & README、许可证、安全说明 \\
30 & 发布与复盘 & Demo、技术文章、下一版本路线 \\
\bottomrule
\end{longtable}

\section{12 周进阶路线}

完成 30 天后，不要立即追逐下一个框架。用三个月深化一个方向：

\begin{itemize}
  \item 第 1--4 周：把毕业项目做到 100 条评测、可观测、可恢复。
  \item 第 5--8 周：选择一个框架重写关键控制流，比较抽象与性能。
  \item 第 9--10 周：阅读 5 篇核心论文，复现一个小实验。
  \item 第 11--12 周：贡献文档、测试或修复到一个上游项目，发布复盘。
\end{itemize}

\section{如何读高速变化的资料}

按四层获取信息：规范和论文用于理解机制；官方文档用于当前接口；成熟开源仓库用于工程模式；社区 Awesome 列表用于发现候选。新闻和短视频适合发现趋势，不适合成为实现依据。

每条资源记录：标题、类型、原始 URL、作者、访问日期、适用阶段、许可证、实践产物和最后核验。季度自动检查链接与仓库活跃度，但是否仍值得推荐需要人工判断。

\section{框架学习策略}

先用手写循环完成一个任务，再选择最匹配的框架。评价框架时看：状态是否显式、恢复是否原生、工具和模型是否可替换、trace 是否完整、部署和版本迁移成本如何。不要同时学五个框架的“Hello World”。

\section{持续复盘模板}

\begin{codeblock}
本周完成的真实任务：
本周新增的评测样本：
最常见失败及根因：
成功率 / P95 延迟 / 单任务成本：
删除了什么不必要复杂度：
下周只改善的一个指标：
\end{codeblock}

\begin{practice}[启动第 1 天]
现在就为一个真实场景完成复杂度评分和任务合同，把它提交到你的项目仓库。明天的任务不是继续收藏资料，而是画出六层数据流。
\end{practice}

\begin{checkpoint}[全书毕业检查]
\begin{itemize}
  \item \checkbox 我能从零写出带校验、状态和终止的 Agent 循环。
  \item \checkbox 我能设计可授权、可恢复、可审计的工具。
  \item \checkbox 我有版本化评测集，能证明改动没有明显回归。
  \item \checkbox 我完成了一个有 Demo、数据和安全说明的项目。
\end{itemize}
\end{checkpoint}

